\kchapter{Setting Up Your Kubernetes Lab}{node}

\begin{objectives}
\begin{itemize}
  \item which local cluster tool to choose (kind, minikube, k3d, Docker Desktop);
  \item how to install \kc{kubectl}, Docker, kind, and Helm on Windows, macOS, and Linux;
  \item how to create a multi-node cluster and verify it works;
  \item what a kubeconfig, a context, and a namespace default are.
\end{itemize}
\end{objectives}

\section{Choosing a local cluster}

You do not need a cloud account to learn Kubernetes. All of these run a real
cluster on your laptop:

\begin{center}
\ktablestart
\begin{tabularx}{\linewidth}{P{2.7cm} L L}
\kheader \kth{Tool} & \kth{How it works} & \kth{Best for} \\
\textbf{kind} & Each node is a Docker container. Multi-node in seconds. &
  CI pipelines, multi-node experiments. \textbf{Used in this book.} \\
\textbf{minikube} & A VM or container with many built-in add-ons. &
  Beginners who like \kc{minikube addons enable \ldots}. \\
\textbf{k3d} & Runs the lightweight k3s distribution in Docker. &
  Very fast, low memory. Built-in load balancer. \\
\textbf{Docker Desktop} & One-click single-node cluster. &
  Simplest possible start on Windows/macOS. \\
\end{tabularx}
\end{center}

\begin{protip}
Give Docker at least \textbf{4 CPUs and 8\,GB RAM} for the later chapters
(especially monitoring). On Windows, use Docker Desktop with the WSL\,2 backend.
\end{protip}

\section{Install the tools}

You need four command-line tools:
\kicon{docker}~\textbf{Docker} (to run containers and kind),
\kicon{k8s-logo}~\textbf{kubectl} (to talk to clusters),
\kicon{k8s-logo}~\textbf{kind} (to create local clusters), and
\kicon{helm}~\textbf{Helm} (the package manager, used from Chapter 18).

\subsection{Windows (PowerShell)}

\begin{bashcode}{PowerShell}
# Install Docker Desktop from https://www.docker.com/products/docker-desktop/
winget install -e --id Kubernetes.kubectl
winget install -e --id Kubernetes.kind
winget install -e --id Helm.Helm
\end{bashcode}

\subsection{macOS (Homebrew)}

\begin{bashcode}{terminal}
brew install --cask docker
brew install kubectl kind helm
\end{bashcode}

\subsection{Linux}

\begin{bashcode}{terminal}
# kubectl (latest stable)
curl -LO "https://dl.k8s.io/release/$(curl -Ls https://dl.k8s.io/release/stable.txt)/bin/linux/amd64/kubectl"
sudo install -m 0755 kubectl /usr/local/bin/kubectl

# kind (check https://kind.sigs.k8s.io for the latest version)
curl -Lo ./kind https://kind.sigs.k8s.io/dl/latest/kind-linux-amd64
sudo install -m 0755 kind /usr/local/bin/kind

# Helm
curl -fsSL https://raw.githubusercontent.com/helm/helm/main/scripts/get-helm-3 | bash
\end{bashcode}

Check everything:

\begin{bashcode}{terminal}
docker version
kubectl version --client
kind version
helm version
\end{bashcode}

\section{Create a multi-node cluster}

A single-node cluster hides interesting behaviour (scheduling, spreading,
node failure). Let us create one control-plane node and two workers.

\begin{yamlcode}{kind-cluster.yaml}
kind: Cluster
apiVersion: kind.x-k8s.io/v1alpha4
name: lab
nodes:
  - role: control-plane
  - role: worker
  - role: worker
\end{yamlcode}

\begin{bashcode}{terminal}
kind create cluster --config kind-cluster.yaml
kubectl cluster-info --context kind-lab
kubectl get nodes -o wide
\end{bashcode}

You should see something like:

\begin{plaincode}{output}
NAME                STATUS   ROLES           AGE   VERSION
lab-control-plane   Ready    control-plane   60s   v1.33.1
lab-worker          Ready    <none>          40s   v1.33.1
lab-worker2         Ready    <none>          40s   v1.33.1
\end{plaincode}

Look inside: the control-plane components themselves run as Pods!

\begin{bashcode}{terminal}
kubectl get pods -n kube-system
# coredns-..., etcd-lab-control-plane, kube-apiserver-lab-control-plane,
# kube-controller-manager-..., kube-scheduler-..., kube-proxy-..., kindnet-...
\end{bashcode}

\section{kubeconfig and contexts}

How does \kc{kubectl} know which cluster to talk to? It reads a
\kterm{kubeconfig} file, by default \kc{\textasciitilde/.kube/config}. That file
holds three lists:

\begin{itemize}
  \item \textbf{clusters} --- API server addresses and CA certificates;
  \item \textbf{users} --- credentials (certificates, tokens, or plugins);
  \item \textbf{contexts} --- a named pairing of \emph{cluster + user + default namespace}.
\end{itemize}

\begin{bashcode}{terminal}
kubectl config get-contexts           # list contexts (* = current)
kubectl config current-context        # e.g. kind-lab
kubectl config use-context kind-lab   # switch cluster
kubectl config set-context --current --namespace=dev   # change default namespace
\end{bashcode}

\begin{warningbox}
The most expensive mistake in Kubernetes is running a command against the
\emph{wrong cluster}. Always check your context before destructive commands.
Tools like \kc{kubectx}/\kc{kubens} and a shell prompt that shows the current
context are worth installing on day one.
\end{warningbox}

\section{Helpful extras}

\begin{itemize}
  \item \textbf{Shell completion}: \kc{kubectl completion bash|zsh|powershell}.
  \item \textbf{An alias}: \kc{alias k=kubectl} (PowerShell:
        \kc{Set-Alias k kubectl}).
  \item \textbf{k9s}: a terminal UI to browse and manage the cluster.
  \item \textbf{stern}: tail logs from many Pods at once.
  \item \textbf{VS Code / Cursor Kubernetes extension}: YAML schema validation
        and auto-completion.
\end{itemize}

\begin{tryit}
Create the \kc{lab} cluster, then delete one worker's container with
\kc{docker stop lab-worker2}. Run \kc{kubectl get nodes -w} and watch the node
turn \kc{NotReady} after about 40 seconds. Start it again with
\kc{docker start lab-worker2}. You just simulated a machine failure.
\end{tryit}

\subsection*{Cleaning up}

\begin{bashcode}{terminal}
kind delete cluster --name lab
\end{bashcode}

\begin{recap}
\begin{itemize}
  \item kind runs a full multi-node cluster inside Docker containers.
  \item \kc{kubectl} reads kubeconfig; a \emph{context} = cluster + user + namespace.
  \item Even control-plane components run as Pods in \kc{kube-system}.
\end{itemize}
\end{recap}
